% Patch record for C-58 -- text-only fixes for debate inconsistencies S-01, S-03, S-04, S-09, S-10.
% Applied 2026-10-28 (Sebas). Source: intake/pending/second_pass_simplification_debate.md (Inconsistency checks).
% No new data: every number used (75.9 ms CI [59.3, 92.6]; 184--534 ms; T_response = 0.0 s) already in results.tex.
% validate_tex.sh OK; compile_pdf.sh OK (102 pages, 0 undefined references).
%
% Diff:
% --- a/sections/methodology.tex
% +++ b/sections/methodology.tex
% @@ -196 +196 @@
% -The architectural framework of the basal ganglia model adopted here is grounded in canonical arbitration circuits widely validated in computational neuroscience for multi-stimulus action selection \citep{Gurney2001a,Gurney2001b,Prescott2006Robot}. Its decentralized winner-take-all dynamics enable low-latency conflict resolution between competing sensory drives. This is evidenced by the near-instantaneous neural decision dispatch measured across experimental trials --- a deterministic $T_{response} = 0.0$~s in simulation and $184$--$534$~ms of neural processing latency across 15 hardware replicates (Section~\ref{sec:results}) --- well within real-time constraints on resource-constrained embedded hardware, without requiring the training cycles or inference infrastructure associated with deep reinforcement learning or model predictive control frameworks, making interpretability and low runtime overhead intrinsic properties of the proposed design.
% +The architectural framework of the basal ganglia model adopted here is grounded in canonical arbitration circuits widely validated in computational neuroscience for multi-stimulus action selection \citep{Gurney2001a,Gurney2001b,Prescott2006Robot}. Its decentralized winner-take-all dynamics enable low-latency conflict resolution between competing sensory drives. \revblue{This is supported by the measured latencies: in simulation, terrain-mode dispatch was deterministic ($T_{response} = 0.0$~s) and stimulus conflicts were resolved with a mean decision latency of $75.9$~ms (Figure~\ref{fig:ablation_combined} and accompanying text); on the physical robot, with the circuits executed on an external computer (Section~\ref{ssec:HardwareSpec}), terrain-driven mode decisions took $184$--$534$~ms across 15 hardware replicates (Figure~\ref{fig:FourPanelLatencies}). The architecture achieves this} without requiring the training cycles or inference infrastructure associated with deep reinforcement learning or model predictive control frameworks, making interpretability and low runtime overhead intrinsic properties of the proposed design.
% --- a/sections/results.tex
% +++ b/sections/results.tex
% @@ -318 +318 @@
% -To quantify the computational efficiency and temporal determinism of the proposed bio-inspired control architecture under stochastic conditions, the neural response latency ($T_{\text{response}}$) and the Empirical Cumulative Distribution Function (ECDF) of the mechanical state-switching delay ($T_{\text{switch}}$) were analyzed across 15 randomized trials for both experimental terrains, the latter illustrated in Figure~\ref{fig:decision_delay}. Both ECDF distributions collapsed into a perfect Heaviside step function: the neural latency was $T_{\text{response}} = 0.0\,\text{s}$ for both terrain types, showing that the threshold-based basal ganglia selection logic filtered out the injected sensor noise and dispatched the correct locomotion command within the same simulation frame as the environmental disturbance. The mechanical transition latency was likewise sub-millisecond in the vast majority of trials, but not fully invariant across terrain types. On rugged terrain, $T_{\text{switch}}$ fell within $0.2$--$0.7\,\text{ms}$ for 13 of 15 trials (a single discrete update cycle of the ROS~2 Control loop operating at $3.3\,\text{kHz}$), while the remaining 2 trials ($13.3\%$) exhibited a long tail out to $\approx 0.6\,\text{s}$, indicating that occasional trials require additional update cycles to complete the locomotion-matrix hot-swap under this terrain's contact conditions. On the inclined slope, $T_{\text{switch}}$ spread across $0.9$--$1.8\,\text{ms}$ in a five-to-six-step staircase rather than collapsing onto a single value, reflecting genuine trial-to-trial dispersion of a comparable order of magnitude rather than a perfectly deterministic delay. \revblue{Neural decision dispatch thus remained deterministic across all 30 trials, while the residual mechanical variability, including the rugged-terrain long tail, is a bounded, terrain-dependent effect rather than chattering.}
% +To quantify the computational efficiency and temporal determinism of the proposed bio-inspired control architecture under stochastic conditions, the neural response latency ($T_{\text{response}}$) and the Empirical Cumulative Distribution Function (ECDF) of the mechanical state-switching delay ($T_{\text{switch}}$) were analyzed across 15 randomized trials for both experimental terrains, the latter illustrated in Figure~\ref{fig:decision_delay}. \revblue{The neural latency was} $T_{\text{response}} = 0.0\,\text{s}$ for both terrain types, showing that the threshold-based basal ganglia selection logic filtered out the injected sensor noise and dispatched the correct locomotion command within the same simulation frame as the environmental disturbance. The mechanical transition latency was likewise sub-millisecond in the vast majority of trials, but not fully invariant across terrain types. On rugged terrain, $T_{\text{switch}}$ fell within $0.2$--$0.7\,\text{ms}$ for 13 of 15 trials (a single discrete update cycle of the ROS~2 Control loop operating at $3.3\,\text{kHz}$), while the remaining 2 trials ($13.3\%$) exhibited a long tail out to $\approx 0.6\,\text{s}$, indicating that occasional trials require additional update cycles to complete the locomotion-matrix hot-swap under this terrain's contact conditions. On the inclined slope, $T_{\text{switch}}$ spread across $0.9$--$1.8\,\text{ms}$ in a five-to-six-step staircase rather than collapsing onto a single value, reflecting genuine trial-to-trial dispersion of a comparable order of magnitude rather than a perfectly deterministic delay. \revblue{Neural decision dispatch thus remained deterministic across all 30 trials, while the residual mechanical variability, including the rugged-terrain long tail, is a bounded, terrain-dependent effect rather than chattering.}
% @@ -515 +515 @@
% -This advantage reverses in the Complex suite, where three stimuli (obstacle, aversive, appetitive) must be handled in sequence: Neural reaches 100\% success while FSM reaches 0\% (all 15 trials timing out or, in one case, tipping over). The mechanism is not that the FSM decides worse while a conflict is active, but that it stays in an unresolved simultaneous conflict (both aversive and appetitive stimuli detected at once) far longer before its fixed priority rule forces a choice: real-time conflict exposure averages $\approx$0.2~s/trial (1.7\% of trial time) for Neural against $\approx$6.1~s/trial (4.7\%) for FSM --- roughly 30$\times$ longer in absolute time. Once the FSM's obstacle-avoidance rule releases control, the robot is frequently repositioned within a second stimulus's detection range, and in the absence of any active rule the controller defaults to an undirected exploratory mode with no residual bias toward the original target, unlike the Neural system's persistent-activation memory (Section~\ref{ssec:FSMBaseline}).
% +This advantage reverses in the Complex suite, where three stimuli (obstacle, aversive, appetitive) must be handled in sequence: Neural reaches 100\% success while FSM reaches 0\% (all 15 trials timing out or, in one case, tipping over). The mechanism is not that the FSM decides worse while a conflict is active, but that it stays in an unresolved simultaneous conflict (both aversive and appetitive stimuli detected at once) far longer before its fixed priority rule forces a choice: real-time conflict exposure averages $\approx$0.2~s/trial (1.7\% of trial time) for Neural against $\approx$6.1~s/trial (4.7\%) for FSM --- roughly 30$\times$ longer in absolute time. Once the FSM's obstacle-avoidance rule releases control, the robot is frequently repositioned within a second stimulus's detection range, and in the absence of any active rule the controller defaults to an undirected exploratory mode with no residual bias toward the original target, unlike the Neural system's persistent-activation memory (Section~\ref{ssec:FSMBaseline}). \revblue{Figure~\ref{fig:mode_timeline_complex} contrasts representative mode traces of both controllers in this suite.}
% @@ -1032 +1032 @@
% -As an open, low-cost exploratory research baseline, the physical robot is constructed using low-torque servomotors and a lightweight PLA structure, which limits its operational envelope compared to industrial-grade platforms. Locomotion gaits are currently executed via pre-configured open-loop control patterns; while this reduces control complexity, it weakens active stabilization capabilities on high-slip surfaces or excessively steep inclinations. \revblue{This constraint stems from the hardware, not from the neural arbitration architecture.} As a representative case, the gait selected during steep-slope negotiation (Figure~\ref{fig:RES_IMU2_real}) reflects the boundaries of the physical actuator envelope: the neural decision module consistently selects the locomotion mode associated with improved stability under the detected terrain conditions, yet the open-loop execution of that gait cannot compensate for the torque deficits inherent to the hardware. Experimental validation was constrained to structured laboratory environments with static stimuli and did not include unconstrained outdoor topographies, dynamic obstacles, or severe external disturbances. Therefore, the empirical results presented herein do not constitute validation of a production-ready system for rugged industrial deployment. Instead, they substantiate the mathematical, biological, and computational feasibility of a highly efficient basal-ganglia inspired arbitration network capable of managing multimodal mobility on resource-constrained embedded hardware.
% +As an open, low-cost exploratory research baseline, the physical robot is constructed using low-torque servomotors and a lightweight PLA structure, which limits its operational envelope compared to industrial-grade platforms. Locomotion gaits are currently executed via pre-configured open-loop control patterns; while this reduces control complexity, it weakens active stabilization capabilities on high-slip surfaces or excessively steep inclinations. \revblue{This constraint stems from the hardware, not from the neural arbitration architecture.} As a representative case, the gait selected during steep-slope negotiation (Figure~\ref{fig:RES_IMU2_real}) reflects the boundaries of the physical actuator envelope: the neural decision module consistently selects the locomotion mode associated with improved stability under the detected terrain conditions, yet the open-loop execution of that gait cannot compensate for the torque deficits inherent to the hardware. Experimental validation was constrained to structured laboratory environments with static stimuli and did not include unconstrained outdoor topographies, dynamic obstacles, or severe external disturbances. Therefore, the empirical results presented herein do not constitute validation of a production-ready system for rugged industrial deployment. Instead, they substantiate the mathematical, biological, and computational feasibility of a highly efficient basal-ganglia inspired arbitration network capable of managing multimodal mobility \revblue{on a low-cost robotic platform}.
% @@ -1052 +1052 @@
% -Beyond the hardware and simulation-environment constraints above, the neural and locomotion-control parameters were manually tuned and kept constant throughout all experiments, without applying adaptive optimization processes that could enhance performance in more diverse scenarios; system resilience was also not assessed under partial hardware failures or temporary sensory dropout conditions likely to occur in real industrial environments. The multimodal sensory integration, moreover, does not follow a formal fusion framework in the classical signal-processing sense---such as Kalman filtering, Bayesian estimation, or deep sensor fusion. Instead, cross-modal integration emerges implicitly within the bio-inspired neural network: each sensory channel feeds differentiated neural pathways that converge through winner-take-all competition dynamics in the basal ganglia module, effectively distributing the fusion function across the network topology. While this approach is biologically grounded, it has not been formally compared against standard fusion methods. Similarly, visual stimulus detection relies on color-based segmentation calibrated under controlled illumination, without evaluation under varying lighting conditions or occlusion. Notwithstanding these design simplifications, the perception-decision noise-robustness evaluation (Section~\ref{ssec:NoiseRobustnessEval}) demonstrated empirical robustness under combined IMU-drift, LiDAR-noise, and camera-illumination-distortion perturbation, evaluated jointly rather than isolated per sensor --- no configuration in this evaluation exercises a single sensor's degradation in isolation from the other two. Within the same evaluation, additive IMU white noise (distinct from the random-walk drift that was exercised) was held at zero throughout and so was not tested here. Robustness was also not uniform across scenarios: the Rugged Terrain suite showed a markedly elevated failure rate under perturbation relative to the other five suites, with 8 tipovers and 4 crashes recorded out of 27 attempts (Section~\ref{sssec:noise_robustness_extra}), indicating that combined sensory perturbation degrades structural stability specifically on this terrain even though the trials that do complete show no corresponding trend toward instability. These constraints do not invalidate the results but indicate that the next development stage should focus on validating the proposed architecture in a prototype that includes robustness and adaptability tests under real, non-controlled conditions.
% +Beyond the hardware and simulation-environment constraints above, the neural and locomotion-control parameters were manually tuned and kept constant throughout all experiments, without applying adaptive optimization processes that could enhance performance in more diverse scenarios; system resilience was also not assessed under partial hardware failures or temporary sensory dropout conditions likely to occur in real industrial environments \revblue{(beyond the simulated camera-detection dropout of Section~\ref{sssec:noise_robustness_extra})}. The multimodal sensory integration, moreover, does not follow a formal fusion framework in the classical signal-processing sense---such as Kalman filtering, Bayesian estimation, or deep sensor fusion. Instead, cross-modal integration emerges implicitly within the bio-inspired neural network: each sensory channel feeds differentiated neural pathways that converge through winner-take-all competition dynamics in the basal ganglia module, effectively distributing the fusion function across the network topology. While this approach is biologically grounded, it has not been formally compared against standard fusion methods. Similarly, visual stimulus detection relies on color-based segmentation calibrated under controlled illumination, without evaluation under varying lighting conditions or occlusion. Notwithstanding these design simplifications, the perception-decision noise-robustness evaluation (Section~\ref{ssec:NoiseRobustnessEval}) demonstrated empirical robustness under combined IMU-drift, LiDAR-noise, and camera-illumination-distortion perturbation, evaluated jointly rather than isolated per sensor --- no configuration in this evaluation exercises a single sensor's degradation in isolation from the other two. Within the same evaluation, additive IMU white noise (distinct from the random-walk drift that was exercised) was held at zero throughout and so was not tested here. Robustness was also not uniform across scenarios: the Rugged Terrain suite showed a markedly elevated failure rate under perturbation relative to the other five suites, with 8 tipovers and 4 crashes recorded out of 27 attempts (Section~\ref{sssec:noise_robustness_extra}), indicating that combined sensory perturbation degrades structural stability specifically on this terrain even though the trials that do complete show no corresponding trend toward instability. These constraints do not invalidate the results but indicate that the next development stage should focus on validating the proposed architecture in a prototype that includes robustness and adaptability tests under real, non-controlled conditions.
% --- a/sections/conclusions.tex
% +++ b/sections/conclusions.tex
% @@ -8 +8 @@
% -In scenarios presenting simultaneous appetitive and aversive stimuli, the basal-ganglia-inspired winner-take-all arbitration module (Section~\ref{sssec:Basal_Ganglia_Module}) consistently prioritized the most relevant objective, its predominant channel activation inhibiting competing channels and reducing the probability of erratic trajectory changes; this behavior yielded direction-switching neural latencies of 184--534~ms (Figure~\ref{fig:FourPanelLatencies}). The alignment between decision-neuron activity peaks and executed maneuvers supports the hypothesis that a basal-ganglia-inspired architecture combined with multimodal sensory processing offers clear advantages in environments where goal prioritization is critical for mission success. 
% +In scenarios presenting simultaneous appetitive and aversive stimuli, the basal-ganglia-inspired winner-take-all arbitration module (Section~\ref{sssec:Basal_Ganglia_Module}) consistently prioritized the most relevant objective, its predominant channel activation inhibiting competing channels and reducing the probability of erratic trajectory changes; \revblue{in simulation, the full circuit resolved such conflicts with a mean decision latency of $75.9$~ms (95\%~CI $[59.3, 92.6]$), while on the physical robot the neural latency of terrain-driven mode transitions was 184--534~ms (Figure~\ref{fig:FourPanelLatencies}).} The alignment between decision-neuron activity peaks and executed maneuvers supports the hypothesis that a basal-ganglia-inspired architecture combined with multimodal sensory processing offers clear advantages in environments where goal prioritization is critical for mission success. 
